\documentclass[10pt,a4paper]{amsart}
\usepackage[margin=19mm]{geometry}
\usepackage{amsmath,amssymb,mathtools}
\usepackage[scr=rsfs,cal=euler]{mathalfa}
\usepackage{xfrac}
\usepackage[unicode,hidelinks,bookmarks=false]{hyperref}
\newcommand{\ie}{i.e.\ }
\newcommand{\cf}{cf.\ }
\newcommand{\resp}{respectively\ }
\newcommand{\Subsection}{\subsection}
\providecommand{\nobreakdash}{\nobreak\hbox{-}}
\setlength{\emergencystretch}{2em}
\multlinegap0pt
\usepackage{amsthm}
\usepackage{enumitem}
\providecommand{\theoremname}{Theorem}
\providecommand{\corollaryname}{Corollary}
\providecommand{\lemmaname}{Lemma}
\providecommand{\remarkname}{Remark}
\providecommand{\definitionname}{Definition}
\providecommand{\examplename}{Example}
\providecommand{\propositionname}{Proposition}
\providecommand{\conjecturename}{Conjecture}
\providecommand{\problemname}{Problem}
\newcommand{\noun}[1]{\textsc{#1}}
\providecommand{\tabularnewline}{\\}
\numberwithin{equation}{section}
\numberwithin{figure}{section}
\theoremstyle{plain}
\newtheorem{thm}{\protect\theoremname}[section]
\theoremstyle{remark}
\newtheorem{rem}[thm]{\protect\remarkname}
\theoremstyle{definition}
\newtheorem{defn}[thm]{\protect\definitionname}
\theoremstyle{plain}
\newtheorem{lem}[thm]{\protect\lemmaname}
\theoremstyle{definition}
\newtheorem{example}[thm]{\protect\examplename}
\theoremstyle{plain}
\newtheorem{cor}[thm]{\protect\corollaryname}
\theoremstyle{plain}
\newtheorem{prop}[thm]{\protect\propositionname}
\theoremstyle{plain}
\newtheorem{conjecture}[thm]{\protect\conjecturename}
\newtheorem{problem}[thm]{\protect\problemname}
\newcommand{\poincare}{Poincar\acute{e}}
\newcommand{\xto}[1]{\xrightarrow{#1}}
\newcommand{\Si}{\Sigma}
\newcommand{\holder}{\rm{H\ddot{o}lder}}
\newcommand{\TIT}{Algebraic Topology}
\newcommand{\NAME}{Homework 3}
\newcommand{\kahler}{\rm{K\ddot{a}hler}}
\newcommand{\kah}{\rm{K\ddot{a}hler}}
\newcommand{\dbar}{\bar{\partial}}
\newcommand{\ddbar}{\partial\bar{\partial}}
\newcommand{\MA}{Monge-Amp\`{e}re}
\newcommand{\KE}{\rm{K\ddot{a}hler–Einstein}}
\newcommand{\pdv}{\partial}
\newcommand{\gar}{\ifmmode{\mathrm{G\mathring{a}rding}}\else\text{Gårding}\fi}
\DeclareMathOperator{\Inv}{\tau}
\newcommand{\Comp}{\operatorname{Comp}}
\newcommand{\SSGF}{S}
\newcommand{\CSGF}{C}
\newcommand{\IGF}{I}
\newcommand{\BGF}{B}
\def\A{\mathcal{A}}
\def\a{\textbf{a}}
\def\b{\textbf{b}}
\def\e{\textbf{e}}
\def\t{\textbf{t}}
\def\E{\mathbb E}
\def\R{\mathbb R}
\def\Z{\mathbb Z}
\def\Q{\mathbb Q}
\def\HH{\mathrm H}
\def\N{\mathbb N}
\def\L{\mathrm L}
\def\P{\mathrm P}
\def\M{\mathrm M}
\def\doubprime{\H}
\def\S{\mathrm S}
\def\G{{\mathrm G}^\mathfrak{a}}
\def\Gh{{\mathrm G}^\mathfrak{h}}
\def\Gah{{\mathrm G}^\mathfrak{ah}}
\def\H {\mathrm H}
\def\GS{\mathrm G}
\def\I{\mathrm I}
\def\U{\mathrm U}
\def\p{\mathbf{p}}
\def\C{\mathcal C}
\def\K{\mathbb K}
\def\x{\mathbf{x}}
\def\y{\mathbf{y}}
\def\u{\mathbf{u}}
\def\v{\mathbf{v}}
\def\w{\mathbf{w}}
\def\z{\mathbf{z}}
\def\MAff{\mathcal {M\!A}}
\def\hom {\mathcal H}
\def\T{\mathrm T}
\def\D{\mathrm D}
\def\RR{\mathrm R}
\def\n{[n]}
\def\proj{\Pi^\downarrow}
\def\pol{\Pi^\uparrow}
\def\Sym{\mathrm {sym}}
\def\CK{\mathbb{C}\mathbb{K}}
\def\V{\mathbb V}
\def\la{\lambda}
\def\La{\Lambda}
\def\ep{\epsilon}
\def\De{\Delta}
\def\de{\delta}
\def\om{\omega}
\def\Om{\Omega}
\def\si{\sigma}
\def\Si{\Sigma}
\def\al{\alpha}
\def\be{\beta}
\def\ga{\gamma}
\def\Ga{\Gamma}
\def\ze{\zeta}
\def\na{\nabla}
\def\<{\langle}
\def\>{\rangle}
\def\sym{\text{sym}}
\def\branden{Br\"and\'en}
\def\vtl{\vartriangleleft}
\providecommand{\corollaryname}{Corollary}
\providecommand{\definitionname}{Definition}
\providecommand{\examplename}{Example}
\providecommand{\factname}{Fact}
\providecommand{\lemmaname}{Lemma}
\providecommand{\propositionname}{Proposition}
\providecommand{\remarkname}{Remark}
\providecommand{\theoremname}{Theorem}
\providecommand{\conjecturename}{Conjecture}
\providecommand{\conjecturename}{Conjecture}
\providecommand{\corollaryname}{Corollary}
\providecommand{\definitionname}{Definition}
\providecommand{\examplename}{Example}
\providecommand{\lemmaname}{Lemma}
\providecommand{\propositionname}{Proposition}
\providecommand{\remarkname}{Remark}
\providecommand{\theoremname}{Theorem}
\providecommand{\problemname}{Problem}
\title[Garding polynomials]{G{\aa}rding Polynomials: the symmetric reduction on the boundary of the Garding cone (Proposition 9.11)}
\author[Hao Fang and Biao Ma]{Hao Fang and Biao Ma}
\date{Source: arXiv:2604.27755v3 (CC BY 4.0)}
\begin{document}
\maketitle
\vspace{1em}
\noindent\textit{Source and licence.} G{\aa}rding Polynomials. Version arXiv:2604.27755v3 (CC BY 4.0), distributed under the Creative Commons Attribution 4.0 International licence, http://creativecommons.org/licenses/by/4.0/, as recorded in the arXiv OAI licence metadata for this exact version and shown on the abstract page https://arxiv.org/abs/2604.27755v3. This item is an editorial re-presentation of the paper's symmetric reduction proposition together with the paper's own proof of it and the source-local chain that proof uses. A full rights notice is delivered at licenses/arxiv-2604.27755-CC-BY-4.0.txt.
\noindent\textit{Editorial scope.} Counted unit: the paper's symmetric reduction on the boundary of the Garding cone, printed as Proposition 8.11 (source0.tex lines 2663--2675), delivered together with the paper's complete original proof of it (source0.tex lines 2677--2728) as the single primary proof body. No mathematics is written, repaired or re-derived: statement and proof are the paper's own text line for line, in the paper's own notation ($\mathcal{B}$, $\mathcal{C}_{\partial_{0}f}$, $\partial_{0}f$, $\u$, $u^{*}$, the interval $I$). Required context retained uncounted, because the counted proof is the paper's own assembly over it: the basic lemma with its proof (source0.tex lines 521--541); a second lemma with its proof that the counted proof invokes (lines 1363--1386); and the two-variable case of the same reduction with its complete proof (lines 2591--2662). Every cross-reference inside the retained spans resolves inside the delivered document, and no citation command occurs in any retained span, so the delivered document carries no bibliography and no reference or citation command of the delivered text was rewritten or adapted. Printed numbers are the paper's own: the wrapper restores the paper's section counter and its single shared section-relative theorem counter before each retained group, so the retained basic lemma prints as Lemma 2.5, the retained second lemma as Lemma 5.15, the retained two-variable case as Proposition 8.10 and the counted statement as Proposition 8.11. Omitted from this package and not redistributed: the paper's preamble and front matter as such (of which only the title and the author names are reproduced in the wrapper); the introduction with the two main theorems; the whole theory of Garding polynomials and their cones outside the retained ranges; the later sections, appendices and examples; and the paper's entire bibliography with its bib data file. Scope note, disclosed: the drawn paper is a long research article whose headline theorems are assembled across many sections, so the counted unit is the paper's symmetric reduction proposition, which is a complete research-level result in its own right and whose proof the paper gives in full in the retained range. Printed slips kept literal, among them the paper's own wording and its own choice of symbols. Layout and class: the paper's own class is the standard AMS amsart class with the 11pt, letterpaper and english options, which is not a publisher class; the delivered wrapper is the lane's pinned amsart editorial wrapper at 10pt on A4, to which this package adds the paper's own theorem-like declarations with their shared section-relative counter, its own {\aa} spelling of the Garding name and its own shorthand macros, all copied verbatim from the paper's source. No publisher class file, no publisher style file and no upstream script is redistributed. Assets: no retained span of this package contains a figure environment or a graphics-inclusion command, no image or artwork file exists in or is redistributed with this package, and the manifest and the delivery evidence both carry an empty asset-dependency list. A full rights notice is delivered at licenses/arxiv-2604.27755-CC-BY-4.0.txt.
\setcounter{section}{2}
\setcounter{thm}{4}
\begin{lem}
\label{lem:basic}Assume $R\subset\R^{n}$.
\begin{enumerate}
\item If $R$ passes PRT, then $R$ is connected; 
\item \label{enu:if--and}if $R_{1},$$R_{2}\in\R^{n}$ both pass PRT, then
$R_{1}\cap R_{2}\neq\emptyset$ ; 
\item \label{enu:directproduct}if $R_{1}\subset\R^{n_{1}}$ and $R_{2}\subset\R^{n_{2}}$
both pass PRT, then $R_{1}\times R_{2}\subset\R^{n_{1}+n_{2}}$ passes
PRT;
\item Fix an index set $I$. If $R_{i}$ passes PRT for all $i\in I$, then
$\cup_{i\in I}R_{i}$ passes PRT. If $\cap_{I}R_{i}\not=\emptyset$
then it also passes PRT.
\item \label{enu:NRTlabel}If $R\subset\R^{n}$ is NRT, so is any of its
open subsets.
\end{enumerate}
\end{lem}

\setcounter{section}{5}
\setcounter{thm}{14}
\begin{lem}
\label{lem:For-any-.}Suppose $f\in\G$. If $\pi_{I}(\mathbf{x}_{0})\in\pi_{I}(\mathcal{C}_{f})$,
then for $\mathbf{y}\in\mathbb{R}^{k}$, 
\[
g(\mathbf{y}):=f(\iota_{I}(\mathbf{y})+\mathbf{x}_{0})\in\G.
\]
In particular, if $\x_{0}\in\C_{f}^{(k_{0})}$ where $k_{0}\leq k$,
then $g(\y)\in\G$.
\end{lem}

\begin{proof}
If $g\equiv c\geq0$, then the result is trivial. Otherwise, denote
\[
R:=\{\mathbf{y}:\iota_{I}(\mathbf{y})+\mathbf{x}_{0}\in\mathcal{C}_{f}\},
\]
which is an open subset of $\{g(\mathbf{y})>0\}$. Since $\pi_{I}(\mathbf{x}_{0})\in\pi_{I}(\mathcal{C}_{f})$,
for $\mathbf{a}=\sum_{i\in I}\mathbf{e}_{i}\in\Gamma_{n,|I|}^{+}$,
there exists $t>>0$ such that $t\mathbf{a}+\mathbf{x}_{0}\in\C_{f}$.
Then $R$ is not empty and passes PRT . Thus, by Lemma \ref{lem:basic},
$R$ is connected. Suppose $\mathbf{y}_{0}\in\pdv R$. Then $\iota_{I}(\mathbf{y}_{0})+\mathbf{x}_{0}\in\pdv\mathcal{C}_{f}$
and hence $g(\mathbf{y}_{0})=0$. Thus, $R$ is both open and closed
in $\{g(\mathbf{y})>0\}$. Hence, $R$ is a connected component of
$\{g(\mathbf{y})>0\}$. Therefore, $g$ is Gårding. 
\end{proof}

\setcounter{section}{8}
\setcounter{thm}{9}
\begin{prop}
\label{prop:n=00003D2}Let $f(x_{0},x_{1},x_{2})\in\mathcal{B}$ be
a multi-affine polynomial and symmetric in $(x_{1},x_{2})$. Let $(u_{0},u_{1},u_{2})\in\partial\mathcal{C}_{\partial_{0}f}$.
Then there exists a diagonal $u$ with $\partial_{0}f(u_{0},u,u)=0$
such that $\min\{u_{1},u_{2}\}\le u\le\max\{u_{1},u_{2}\},$ and 
\[
f(u_{0},u_{1},u_{2})\le f(u_{0},u,u).
\]
Moreover, if $u_{1}\neq u_{2}$ then $\min\{u_{1},u_{2}\}<u<\max\{u_{1},u_{2}\}.$
\end{prop}

\begin{proof}
Considering symmetry, write
\[
f(x_{0},x_{1},x_{2})=ax_{0}x_{1}x_{2}+b_{0}x_{1}x_{2}+bx_{0}(x_{1}+x_{2})+c_{0}x_{0}+c(x_{1}+x_{2})+d.
\]
Then 
\[
\partial_{0}f(x_{1},x_{2})=ax_{1}x_{2}+b(x_{1}+x_{2})+c_{0},\qquad f=x_{0}\,\partial_{0}f+h,
\]
where $h(x_{1},x_{2})=b_{0}x_{1}x_{2}+c(x_{1}+x_{2})+d$.

Since $f\in\mathcal{B}$ polynomials $\partial_{0}f$ and $\partial_{1}f$
are G\aa rding with degrees at most 2 and hence real--stable. In
particular, 
\[
b^{2}-ac_{0}\ge0\qquad\text{and}\qquad bb_{0}-ac\ge0.
\]

Define $p(x)=\partial_{0}f(x,x)=ax^{2}+2bx+c_{0}$. Since $p$ is
stable, its G\aa rding component is the interval to the right of
its largest real root. Let $u$ be this boundary point, i.e. 
\[
u=\frac{-b+\sqrt{\,b^{2}-ac_{0}\,}}{a}\quad(a\ne0),\ \text{or }u=-\frac{c_{0}}{2b}\quad(a=0).
\]
Then $\partial_{0}f(u,u)=0$, so $(u_{0},u,u)\in\partial\mathcal{C}_{\partial_{0}f}$.
Because $(u_{0},u_{1},u_{2})\in\partial\mathcal{C}_{\partial_{0}f}$,
we also have $\partial_{0}f(u_{1},u_{2})=0$. Using $f=x_{0}\partial_{0}f+h$,
\[
f(u_{0},u_{1},u_{2})-f(u_{0},u,u)=h(u_{1},u_{2})-h(u,u).
\]

If $a=0$, then $\partial_{0}f=b(x_{1}+x_{2})+c_{0}$ and $\partial_{0}f(u_{1},u_{2})=0$
imply  $u=(u_{1}+u_{2})/2$. Hence $\min\{u_{1},u_{2}\}<u<\max\{u_{1},u_{2}\}$
when $u_{1}\ne u_{2}$. A direct computation gives 
\[
h(u_{1},u_{2})-h(u,u)=-\frac{b_{0}}{4}(u_{1}-u_{2})^{2}\le0,
\]
and since $\partial_{1}f$ is G\aa rding we have $b_{0}\ge0$. Thus
$f(u_{0},u_{1},u_{2})\le f(u_{0},u,u)$.

Now assume $a>0$. Set $A_{1}:=au_{1}+b$, $A_{2}:=au_{2}+b$, and
$A:=au+b$. Then $A_{1}=\partial_{02}f(u_{1},u_{2})$, $A_{2}=\partial_{01}f(u_{1},u_{2})$,
and $A=\partial_{01}f(u,u)$, all strictly positive by boundary non-degeneracy.
From $\partial_{0}f(u_{1},u_{2})=0$ one finds 
\[
A_{1}A_{2}=b^{2}-ac_{0},
\]
while by definition $A=\sqrt{b^{2}-ac_{0}}$, hence $A^{2}=A_{1}A_{2}$.
If $u_{1}\ne u_{2}$, then $A_{1}\ne A_{2}$; assuming that $A_{1}<A_{2}$
gives $A^{2}=A_{1}A_{2}<A_{2}^{2}$, so $A<A_{2}$. Since $x\mapsto ax+b$
is strictly increasing, this implies $u<u_{2}$, and symmetrically
$u>u_{1}$.

Eliminating products using $\partial_{0}f(u_{1},u_{2})=\partial_{0}f(u,u)=0$
yields 
\[
h(u_{1},u_{2})-h(u,u)=-\frac{bb_{0}-ac}{a}\,(u_{1}+u_{2}-2u).
\]
Because $bb_{0}-ac\ge0$ and $u_{1}+u_{2}-2u\ge0$, the right--hand
side is nonpositive. Hence $f(u_{0},u_{1},u_{2})\le f(u_{0},u,u)$. 
\end{proof}

\begin{prop}
\label{prop:symmetry}Let $f\in\mathcal{B}$ be multi-affine and symmetric
in $x_{1},\dots,x_{n}$. Let $(u_{0},\u)=(u_{0},u_{1},\dots,u_{n})\in\partial\mathcal{C}_{\partial_{0}f}$
and set $I=[\min_{i}u_{i},\max_{i}u_{i}]$. Then there exists $u^{*}\in I$
such that 
\[
(u_{0},\u^{*})=(u_{0},u^{*},\dots,u^{*})\in\partial\mathcal{C}_{\partial_{0}f}
\]
and 
\[
f(u_{0},u_{1},\dots,u_{n})\le f(u_{0},u^{*},\dots,u^{*}).
\]
\end{prop}

\begin{proof}
Define $\phi(\x)=\max_{i\in[n]}x_{i}-\min_{i\in[n]}x_{i}$ for $\x\in I^{n}$,
and set 
\[
\mathcal{F}=\bigl\{\x\in I^{n}:\ (u_{0},\x)\in\partial\mathcal{C}_{\partial_{0}f}\ \text{and}\ f(u_{0},\x)\ge f(u_{0},u_{1},\dots,u_{n})\bigr\}.
\]
The set $\mathcal{F}\not=\emptyset$ and is compact, hence, $\phi$ attains
a minimum at some $\mathbf{v}=(v_{1},\dots,v_{n})\in\mathcal{F}$.

By the non-degeneracy condition in the definition of $\mathcal{B}$,
there exists $T\ge0$ such that 
\[
\partial_{i}f(u_{0}+T,\mathbf{v})>0\qquad\text{for all }i=1,\dots,n.
\]
Since $\partial_{0}f\equiv0$ on $\mathcal{F}$, $f(u_{0},\v)=f(u_{0}+T,\v)$
and replacing $u_{0}$ by $u_{0}+T$ does not affect $\mathcal{F}$
or $\phi$, and $\mathbf{v}$ remains a minimizer. 

Suppose that not all coordinates of $\v$ coincide. Let $w^{(1)}<\cdots<w^{(m)}$
be the distinct values among the set $\{v_{i}\}_{i\in[n]}$ with $m\ge2$,
and let 
\[
K=\{k:\ v_{k}=w^{(m)}\}.
\]
Fix an index $i$ with $v_{i}=w^{(1)}$. For each $j\in K$, freeze
every coordinate $x_{k}$ with $k\notin\{i,j\}$ to $v_{k}$ and consider
the three--variable restriction 
\[
g_{ij}(x_0,x_{i},x_{j})=f(x_0,\dots,x_{i},\dots,x_{j},\dots).
\]
Because $\partial_{\ell}f(u_{0}+T,\mathbf{v})>0$ for all frozen indices
$\ell$, the polynomial $g_{ij}$ lies in $\mathcal{B}$ by Lemma \ref{lem:For-any-.} and Proposition
\ref{prop:n=00003D2} applies and we may replace $(v_{i},v_{j})$
by $(u_{ij},u_{ij})$ with 
\[
u_{ij}\in[\min\{v_{i},v_{j}\},\max\{v_{i},v_{j}\}),
\]
hence $u_{ij}\neq w^{(m)}$. Thus, each such replacement strictly decreases
$|K|$ which is the multiplicity of the maximal value $w^{(m)}$.
After at most $|K|$ steps, the value $w^{(m)}$ disappears from the
coordinate list, so the maximum coordinate strictly decreases and
$\phi$ strictly decreases. This contradicts the minimality of $\phi(\mathbf{v})$.

Therefore, $\v=(u^{*},\dots,u^{*})$ for some $u^{*}\in I$. Since
$\v\in\mathcal{F}$, we have $(u_{0}+T,\u^{*})\in\partial\mathcal{C}_{\partial_{0}f}$.
Hence $(u_{0},\u^{*})\in\pdv\C_{\pdv_{0}f}$ and $f(u_{0},\u^{*})=f(u_{0}+T,\u^{*})$.
By the definition of $\mathcal{F}$, 
\[
f(u_{0},u_{1},\dots,u_{n})\le f(u_{0}+T,u^{*},\dots,u^{*})=f(u_{0},\u^{*}).
\]
We have finished the proof.
\end{proof}

\end{document}
